% ============================================================
% Lampiran: Kode Pengujian
% Simpan sebagai, misalnya: src/chapters/ta/lampiran-pengujian.tex
% Daftarkan di dokumen utama dengan \appendix lalu \input{...}
% ============================================================

\chapter{Kode Pengujian}
\label{app:kode-pengujian}

Lampiran ini memuat seluruh kode yang digunakan dalam pengujian nonfungsional
(P-KNF-01 hingga P-KNF-08) dan pengujian integrasi (P-INT-01). Kode diorganisasi
mulai dari fungsi pendukung, implementasi fungsi verifikasi yang diuji, kemudian
setiap kasus pengujian secara berurutan. Seluruh kode berada pada berkas
\texttt{audit-trail/src/tests.rs}, kecuali \texttt{verify\_record} dan
\texttt{verify\_chain} yang berada pada \texttt{audit-trail/src/utils.rs}.

% ─────────────────────────────────────────────────────────────
\section{Fungsi Pendukung Pengujian}
% ─────────────────────────────────────────────────────────────

\begin{listing}[ht]
\begin{minted}[linenos, breaklines]{rust}
fn make_ev4_event() -> AuditEvent {
    AuditEvent {
        source_component: AuditSourceComponent::HospitalClient,
        source_timestamp: Utc::now(),
        event: Event {
            actor_id: "0xabc123_tenaga_medis".to_string(),
            actor_type: AuditActorType::MedicalPersonnel,
            target_object_type: AuditTargetObjectType::MedicalRecord,
            target_object: "ipfs://bafybeig_cid_rme".to_string(),
            outcome: AuditOutcome::Success,
            action_type: AuditActionType::Read,
            details: AuditEventDetails::MedicalRecordAccess {
                patient_iota_address: "0xpatient_iota_addr".to_string(),
                record_index: Some(0),
                role_used: "MedicalPersonnel".to_string(),
            },
        },
    }
}

fn make_ev1_event(actor: &str, outcome: AuditOutcome) -> AuditEvent {
    AuditEvent {
        source_component: AuditSourceComponent::HospitalClient,
        source_timestamp: Utc::now(),
        event: Event {
            actor_id: actor.to_string(),
            actor_type: AuditActorType::MedicalPersonnel,
            target_object_type: AuditTargetObjectType::AccessKeys,
            target_object: "session".to_string(),
            outcome,
            action_type: AuditActionType::SignIn,
            details: AuditEventDetails::Authentication {
                auth_method: "iota_signature".to_string(),
                role: "MedicalPersonnel".to_string(),
                failure_reason: None,
            },
        },
    }
}
\end{minted}
\caption{Fungsi pembantu pembangkit \texttt{AuditEvent} uji}
\label{code:helper-make-event}
\end{listing}

\begin{listing}[ht]
\begin{minted}[linenos, breaklines]{rust}
fn build_chain(n: usize) -> Vec<AuditRecord> {
    let mut records = Vec::new();
    let mut prev_hash: Option<String> = None;

    for i in 0..n {
        let event = make_ev1_event(
            &format!("actor_{i}"),
            AuditOutcome::Success,
        );
        let record = create_audit_record(event, prev_hash.clone());
        prev_hash = Some(record.record_hash.clone());
        records.push(record);
    }

    records
}

fn make_signed_event(event: &AuditEvent) -> (SignedEvent, SigningKey) {
    let signing_key = SigningKey::generate(&mut OsRng);
    let payload = serde_json::to_string(event).expect("serialize event");
    let sig = signing_key.sign(payload.as_bytes());

    let signed = SignedEvent {
        payload,
        signature: hex::encode(sig.to_bytes()),
        public_key: hex::encode(signing_key.verifying_key().as_bytes()),
    };

    (signed, signing_key)
}
\end{minted}
\caption{Fungsi pembantu pembangkit rantai \textit{record} dan \texttt{SignedEvent}}
\label{code:helper-chain-signed}
\end{listing}


\begin{listing}[ht]
\begin{minted}[linenos, breaklines]{rust}
// Ditambahkan ke dalam impl Utils di audit-trail/src/utils.rs

pub fn verify_record(record: &AuditRecord) -> bool {
    let expected = Self::calculate_record_hash(
        record.record_id,
        record.timestamp,
        record.prev_record_hash.clone(),
        &record.event,
    );
    record.record_hash == expected
}

pub fn verify_chain(records: &[AuditRecord]) -> Result<(), String> {
    let mut prev_hash: Option<String> = None;

    for (i, record) in records.iter().enumerate() {
        if record.prev_record_hash != prev_hash {
            return Err(format!(
                "record {}: prev_hash mismatch (expected {:?}, got {:?})",
                i + 1,
                prev_hash,
                record.prev_record_hash
            ));
        }

        if !Self::verify_record(record) {
            return Err(format!("record {}: hash mismatch", i + 1));
        }

        prev_hash = Some(record.record_hash.clone());
    }

    Ok(())
}
\end{minted}
\caption{Implementasi \texttt{verify\_record} dan \texttt{verify\_chain} pada \texttt{Utils}}
\label{code:impl-verify}
\end{listing}

\begin{listing}[ht]
\begin{minted}[linenos, breaklines]{rust}
#[test]
fn verify_record_valid() {
    let event = make_ev4_event();
    let record = create_audit_record(event, None);
    assert!(
        Utils::verify_record(&record),
        "Record yang baru dibuat harus valid"
    );
}
\end{minted}
\caption{Pengujian P-KNF-01a: \textit{record} valid diterima}
\label{code:p-knf-01a}
\end{listing}

\begin{listing}[ht]
\begin{minted}[linenos, breaklines]{rust}
#[test]
fn verify_record_tampered_actor_id() {
    let event = make_ev4_event();
    let mut record = create_audit_record(event, None);

    // Modifikasi: ganti actor_id dengan nilai palsu
    record.event.event.actor_id = "0xpenyerang_palsu".to_string();

    assert!(
        !Utils::verify_record(&record),
        "Record yang dimodifikasi harus gagal verifikasi"
    );
}

#[test]
fn verify_record_tampered_outcome() {
    let event = make_ev4_event();
    let mut record = create_audit_record(event, None);

    // Modifikasi: ubah outcome dari Success menjadi Failure
    record.event.event.outcome = AuditOutcome::Failure;

    assert!(
        !Utils::verify_record(&record),
        "Perubahan outcome harus terdeteksi"
    );
}
\end{minted}
\caption{Pengujian P-KNF-01b: modifikasi \textit{record} terdeteksi}
\label{code:p-knf-01b}
\end{listing}


\begin{listing}[ht]
\begin{minted}[linenos, breaklines]{rust}
#[test]
fn verify_chain_three_records_valid() {
    let records = build_chain(3);
    assert!(
        Utils::verify_chain(&records).is_ok(),
        "Rantai tiga record yang tidak dimodifikasi harus valid"
    );
}

#[test]
fn verify_chain_single_record_valid() {
    let records = build_chain(1);
    assert!(Utils::verify_chain(&records).is_ok());
}

#[test]
fn verify_chain_empty() {
    let records: Vec<AuditRecord> = vec![];
    assert!(Utils::verify_chain(&records).is_ok());
}
\end{minted}
\caption{Pengujian P-KNF-02a: rantai \textit{hash} valid}
\label{code:p-knf-02a}
\end{listing}

\begin{listing}[ht]
\begin{minted}[linenos, breaklines]{rust}
#[test]
fn verify_chain_tamper_middle_record() {
    let mut records = build_chain(3);

    // Modifikasi record ke-2: ubah actor_id
    records[1].event.event.actor_id = "0xpenyerang".to_string();

    let result = Utils::verify_chain(&records);
    assert!(
        result.is_err(),
        "Perusakan pada record ke-2 harus terdeteksi"
    );

    let msg = result.unwrap_err();
    assert!(
        msg.contains("record 2"),
        "Pesan error harus menyebut record 2, dapat: {msg}"
    );
}

#[test]
fn verify_chain_tamper_first_record() {
    let mut records = build_chain(3);

    // Modifikasi record pertama
    records[0].event.event.actor_id = "0xpenyerang".to_string();

    let result = Utils::verify_chain(&records);
    assert!(result.is_err(), "Perusakan record ke-1 harus terdeteksi");

    let msg = result.unwrap_err();
    assert!(
        msg.contains("record 1") || msg.contains("record 2"),
        "Efek domino harus terdeteksi: {msg}"
    );
}
\end{minted}
\caption{Pengujian P-KNF-02b: perusakan rantai terdeteksi}
\label{code:p-knf-02b}
\end{listing}


\begin{listing}[ht]
\begin{minted}[linenos, breaklines]{rust}
#[test]
fn format_roundtrip_single_record() {
    let event = make_ev4_event();
    let record = create_audit_record(event, None);

    // Serialisasi ke JSON (simulasi satu baris JSON-Lines)
    let json_line = serde_json::to_string(&record)
        .expect("serialize record");

    // Deserialisasi kembali ke AuditRecord
    let decoded: AuditRecord = serde_json::from_str(&json_line)
        .expect("deserialize record");

    assert_eq!(record.record_id, decoded.record_id);
    assert_eq!(record.record_hash, decoded.record_hash);
    assert_eq!(
        record.event.event.actor_id,
        decoded.event.event.actor_id
    );
}
\end{minted}
\caption{Pengujian P-KNF-03a: \textit{round-trip} satu \textit{record} JSON-Lines}
\label{code:p-knf-03a}
\end{listing}

\begin{listing}[ht]
\begin{minted}[linenos, breaklines]{rust}
#[test]
fn format_roundtrip_chain_five_records() {
    let records = build_chain(5);

    // Simulasi penulisan berkas log (JSON-Lines)
    let jsonl: String = records
        .iter()
        .map(|r| serde_json::to_string(r).unwrap())
        .collect::<Vec<_>>()
        .join("\n");

    // Pembacaan kembali baris per baris
    let decoded: Vec<AuditRecord> = jsonl
        .lines()
        .map(|line| serde_json::from_str(line)
            .expect("tiap baris harus valid JSON"))
        .collect();

    assert_eq!(records.len(), decoded.len(),
        "Jumlah record harus sama");

    for (original, parsed) in records.iter().zip(decoded.iter()) {
        assert_eq!(original.record_id, parsed.record_id);
        assert_eq!(original.record_hash, parsed.record_hash);
        assert_eq!(original.prev_record_hash, parsed.prev_record_hash);
    }

    // Rantai yang dibaca kembali harus tetap valid
    assert!(Utils::verify_chain(&decoded).is_ok());
}
\end{minted}
\caption{Pengujian P-KNF-03b: \textit{round-trip} lima \textit{record} dan verifikasi rantai}
\label{code:p-knf-03b}
\end{listing}


\begin{listing}[ht]
\begin{minted}[linenos, breaklines]{rust}
#[test]
fn signed_event_valid_accepted() {
    let event = make_ev1_event("0xactor_sah", AuditOutcome::Success);
    let (signed, _key) = make_signed_event(&event);

    let result = Utils::verify_and_extract_event(&signed);
    assert!(
        result.is_ok(),
        "Event dengan tanda tangan valid harus diterima: {:?}",
        result.err()
    );

    let extracted = result.unwrap();
    assert_eq!(extracted.event.actor_id, "0xactor_sah");
}
\end{minted}
\caption{Pengujian P-KNF-06: \texttt{SignedEvent} valid diterima}
\label{code:p-knf-06}
\end{listing}

\begin{listing}[ht]
\begin{minted}[linenos, breaklines]{rust}
#[test]
fn signed_event_tampered_payload_rejected() {
    let event = make_ev1_event("0xactor_sah", AuditOutcome::Success);
    let (mut signed, _key) = make_signed_event(&event);

    // Modifikasi payload setelah penandatanganan
    signed.payload.push_str("x");

    let result = Utils::verify_and_extract_event(&signed);
    assert!(
        result.is_err(),
        "Payload yang dimodifikasi harus ditolak"
    );

    let err = result.unwrap_err().to_string();
    assert!(
        err.to_lowercase().contains("signature")
            || err.to_lowercase().contains("tanda tangan")
            || err.to_lowercase().contains("invalid"),
        "Pesan error harus mengindikasikan kegagalan verifikasi: {err}"
    );
}
\end{minted}
\caption{Pengujian P-KNF-07: \textit{payload} yang dimodifikasi ditolak}
\label{code:p-knf-07}
\end{listing}



\begin{listing}[ht]
\begin{minted}[linenos, breaklines]{rust}
#[test]
fn signed_event_wrong_key_rejected() {
    let event = make_ev1_event("0xactor_sah", AuditOutcome::Success);
    let (mut signed, _original_key) = make_signed_event(&event);

    // Ganti signature dengan tanda tangan dari keypair yang berbeda
    let different_key = SigningKey::generate(&mut OsRng);
    let wrong_sig = different_key.sign(signed.payload.as_bytes());
    signed.signature = hex::encode(wrong_sig.to_bytes());
    // public_key masih milik original_key → mismatch

    let result = Utils::verify_and_extract_event(&signed);
    assert!(
        result.is_err(),
        "Tanda tangan dari kunci berbeda harus ditolak"
    );
}

#[test]
fn signed_event_wrong_public_key_rejected() {
    let event = make_ev1_event("0xactor_sah", AuditOutcome::Success);
    let (mut signed, _key) = make_signed_event(&event);

    // Ganti public_key dengan kunci publik acak yang tidak bersesuaian
    let impostor_key = SigningKey::generate(&mut OsRng);
    signed.public_key =
        hex::encode(impostor_key.verifying_key().as_bytes());

    let result = Utils::verify_and_extract_event(&signed);
    assert!(
        result.is_err(),
        "Public key yang tidak bersesuaian harus ditolak"
    );
}
\end{minted}
\caption{Pengujian P-KNF-08: tanda tangan dari kunci yang tidak bersesuaian ditolak}
\label{code:p-knf-08}
\end{listing}

\begin{listing}[ht]
\begin{minted}[linenos, breaklines]{rust}
#[test]
fn full_cycle_ev4_chain_valid() {
    // Langkah 1: Bangkitkan event EV4 (Medical Record Access)
    let ev4 = make_ev4_event();

    // Langkah 2: Bangun tiga record berurutan; record ke-2 adalah EV4
    let rec1 = create_audit_record(
        make_ev1_event("actor_1", AuditOutcome::Success),
        None,
    );
    let rec2 = create_audit_record(
        ev4,
        Some(rec1.record_hash.clone()),
    );
    let rec3 = create_audit_record(
        make_ev1_event("actor_3", AuditOutcome::Success),
        Some(rec2.record_hash.clone()),
    );

    // Langkah 3: Verifikasi hubungan prev_record_hash antar-record
    assert_eq!(
        rec2.prev_record_hash,
        Some(rec1.record_hash.clone()),
        "prev_record_hash rec2 harus = record_hash rec1"
    );
    assert_eq!(
        rec3.prev_record_hash,
        Some(rec2.record_hash.clone()),
        "prev_record_hash rec3 harus = record_hash rec2"
    );
\end{minted}
\caption{Pengujian P-INT-01 (bagian 1): pembangkitan \textit{event} EV4 dan verifikasi hubungan rantai}
\label{code:p-int-01a}
\end{listing}

\begin{listing}[ht]
\begin{minted}[linenos, breaklines]{rust}
    // Langkah 4a: Rantai valid sebelum modifikasi
    let chain = vec![rec1.clone(), rec2.clone(), rec3.clone()];
    assert!(
        Utils::verify_chain(&chain).is_ok(),
        "Rantai tiga record harus valid sebelum modifikasi"
    );

    // Langkah 4b: Rusak record ke-2, periksa efek domino
    let mut tampered = chain.clone();
    tampered[1].event.event.actor_id = "0xpenyerang_palsu".to_string();

    let result = Utils::verify_chain(&tampered);
    assert!(result.is_err(),
        "Perusakan pada record ke-2 harus terdeteksi");

    // Langkah 5: Bukti minimal dari record EV4
    let ev4_record = &chain[1];
    assert_eq!(
        ev4_record.event.event.actor_id,
        "0xabc123_tenaga_medis"
    );
    assert!(matches!(
        ev4_record.event.event.outcome,
        AuditOutcome::Success
    ));

    // Langkah 6: Format JSON-Lines round-trip tetap valid
    let jsonl = chain.iter()
        .map(|r| serde_json::to_string(r).unwrap())
        .collect::<Vec<_>>()
        .join("\n");

    let decoded: Vec<AuditRecord> = jsonl.lines()
        .map(|l| serde_json::from_str(l).unwrap())
        .collect();

    assert!(Utils::verify_chain(&decoded).is_ok());
}
\end{minted}
\caption{Pengujian P-INT-01 (bagian 2): verifikasi integritas, deteksi perusakan, dan \textit{round-trip} JSON-Lines}
\label{code:p-int-01b}
\end{listing}


\begin{listing}[ht]
\begin{minted}[linenos, breaklines]{rust}
#[test]
fn first_record_has_null_prev_hash() {
    let event = make_ev4_event();
    let record = create_audit_record(event, None);
    assert_eq!(record.prev_record_hash, None,
        "Record pertama dalam rantai harus punya prev_hash null");
}

#[test]
fn record_hash_is_64_char_hex() {
    let event = make_ev4_event();
    let record = create_audit_record(event, None);

    assert_eq!(
        record.record_hash.len(),
        64,
        "SHA-256 hex harus tepat 64 karakter"
    );
    assert!(
        record.record_hash.chars().all(|c| c.is_ascii_hexdigit()),
        "record_hash harus berupa karakter heksadesimal saja"
    );
}
\end{minted}
\caption{Pengujian \textit{sanity check}: format \texttt{prev\_record\_hash} dan \texttt{record\_hash}}
\label{code:sanity}
\end{listing}